% =====================================================================
%   附录 F：横向思维谜题 (Lateral Thinking Puzzles)
% =====================================================================

\section{横向思维谜题}\label{app:ltp}

\subsection{数据集细节}

\subsubsection*{构造细节}

每条样本由一对故事（谜面，例如：\emph{"一个男子走进一家餐厅，点了一碗
甲鱼汤，喝完后便自杀了。他为什么这么做？"}）与真相（谜底）构成。
我们将样本按难度分为四级：简单、中等、困难与专家级。LLM 智能体参与的
LTP 规则如下：
\begin{itemize}
    \item \textbf{角色}：LTP 评测中的角色为主持人与解题者。主持人知晓故事
          与真相，将故事呈现给解题者，并引导其猜出真相。解题者由 LLM
          扮演，通过提出问题并综合主持人的回答来尝试破解真相。
    \item \textbf{解题步骤}：每局游戏设有最大轮数（如 25）。解题者每轮需
          基于已知事实提出一个问题。问题应能以"是"、"否"或"无关"作答。
          主持人以正确答案回应。为降低 LLM 智能体的难度，主持人在解题者
          陷入错误推理方向时，会在回答中提供一定的提示。
    \item \textbf{游戏终止}：当解题者认为自己已猜出真相的主要部分时，可
          向主持人宣告自己的推测剧情。若正确，主持人将宣布游戏结束。
\end{itemize}

\subsubsection*{评测设置}

对每对故事与真相，我们按以下步骤评估模型：
\begin{enumerate}
    \item \textbf{初始化}：通过本地 Python 包安装或 Web API 设置 LTP
          主持人系统。
    \item \textbf{交互}：我们为 LLM 配置系统提示以构建其玩家角色。LLM 在
          每局游戏的最大轮数内担任解题者，只要 LLM 不超出最大 token 长度，
          就可持续进行。在自动评测中，我们将答案限制为"是"、"否"或"无关"，
          并从 \code{gpt-3.5-turbo} 的回复中提取答案。同时，我们要求 LLM
          在自动评测中总结其推理过程，以便终止检测更为准确。
    \item \textbf{校验}：我们对每个 LLM 进行预实验，收集游戏进程中的所有
          情形并设计校验方案。在自动评测中，我们为 \code{gpt-3.5-turbo}
          设置一些关键词以提醒其考虑同义词等灵活情形。
\end{enumerate}

\subsubsection*{评测指标}

我们通过以下自创指标来评估 LLM 的横向推理能力：
\begin{itemize}
    \item \textbf{单局准确率（Single Game Accuracy, SGA）}：在单局游戏中，
          LLM 接近真相的轮次占比。
    \item \textbf{轮次效率（Round Efficiency, RE）}：模型在最大轮数内猜出
          真相的速度。
    \item \textbf{问题相关性（Query Relevance, QR）}：模型所提问题与真相之间
          的相关程度。
    \item \textbf{游戏进度（Game Progress, GP）}：游戏结束前所取得的进度，
          作为主要评测指标。我们将金标准真相分解为若干关键节点，并度量
          智能体所命中的关键节点数量。
\end{itemize}

\subsection{LTP 系统的评测}

我们通过人工验证的方式评估 LTP 系统的准确率，验证系统在关键节点识别与
事实核查方面的准确性。我们比较了自动评测与人工评测下的单局准确率与
问题相关性，发现自动评测有时对智能体更为宽容——尤其在开源模型上——
使得 SGA 与 QR 的得分看上去优于人工评测。我们计划专门训练一个模型作为
游戏的主持人，以提供更佳的游戏体验和更精确的评测。对于游戏进度与
轮次效率，LTP 系统能够提供客观的评测，其水平可与人工评测相匹配。

\subsection{LTP 游戏进度与终止}

游戏进度定义为关键节点被命中的比例。关键节点由 \code{gpt-3.5-turbo}
进行总结，并作为 \code{answer_keys} 字段纳入数据集（示例如下）：
\begin{quote}\small
\textbf{真相（Truth）}：\emph{那天晚上，他们去废弃建筑数台阶。他们在网上
看到传言，称夜晚数台阶会比实际少一。然而第二天再去验证时，发现根本没有
台阶。}\\
\\
\textbf{关键节点（Key points）}：
\begin{enumerate}
    \item 他们想数清废弃建筑的台阶数。
    \item 发生了超自然事件。
    \item 他们在网上看到传言：夜晚数台阶会比实际少一。
    \item 第二天去废弃建筑验证时，发现根本没有台阶。
    \item 他们因极度恐惧而崩溃。
\end{enumerate}
\end{quote}

不同样本的关键节点数量各异。关于是否认定智能体猜中关键节点的判定，我们
首先将相关问题转化为陈述句，再将句子简化为一句。若智能体已猜中某一关键
节点，则删除该节点及相关推理，以避免重复猜测。

\subsection{提示词示例}

\textbf{智能体提示词}如下：
\begin{quote}\small\ttfamily
You are a game player, and you are playing Lateral Thinking Puzzle, also
known as Situation Puzzle. Lateral Thinking Puzzle is a deductive reasoning
game, and here are the game rules:\\
\\
1. At the beginning of the game, you will receive a narrative, referred
to as ``story''. Based on the story, you need to ask questions that can
be answered with ``yes'', ``no'', or ``irrelevant'' to guess out the
``truth''.\\
2. By asking questions, you narrow down the range of possibilities until
you eventually guess out the truth.\\
3. Each time, you can only ask one question.\\
4. Remember that your role is a player. You cannot declare the end of the
game, give up on reasoning, or request a new game.\\
5. You cannot directly repeat information already provided in the story.\\
6. You cannot directly ask for details about the story in the form of
``why'' questions; you need to make your own guesses for truth.\\
7. You cannot directly inquire about the story; you must make your own
deductions.\\
\\
Next, please make full use of the information provided above to engage in
game reasoning. Keep in mind that your questions should be answerable
with ``yes'', ``no'', or ``irrelevant'', and you can only ask one
question at a time.\\
\\
Here is your story: \code{\{story\}}\\
\\
You can start guessing the content of the truth, and I will answer your
questions. Please note that your questions should be answerable with
``yes'', ``no'', or ``irrelevant''.
\end{quote}

\textbf{主持人提示词}：主持人侧的提示词主要包含故事、真相、若干条严格的
游戏规则（参见原文）。其完整结构与上述智能体提示词类似，区别在于主持人
知晓真相，并能根据规则对解题者的提问以"是"、"否"或"无关"作答，并在
特定情形下提供适度的引导性提示。
